% 03_methods.tex — Methods (emulator, oracle, F-and-S-and-M triad, metrics).
% Drafted by item P2-E8-4. Voice: ../STYLE.md (Andreas Maier). Numbers: committed
% leaderboard/faithful_demo + experiment_design.md §0–§3. Figure: \ref{fig:platform}.

\section{Methods}\label{sec:methods}

Our method rests on a single idea: we study an interpretability question on a machine whose
answer we can compute exactly. We describe in turn the substrate that makes the ground truth
obtainable, the oracle that delivers it, the tiered semantics we attach to it, and the
correctness triad that turns ``do we understand?'' into a set of numbers. The per-method
setups live in the Supplementary Information; Fig.~\ref{fig:platform} sketches the pipeline.

\begin{figure}[t]
  \centering
  \includegraphics[width=\linewidth]{fig1_platform_oracle.pdf}
  \caption{\textbf{The measurement platform and its ground-truth oracle.} A
  differentiable, bit-exact port of the Atari VCS (left), realized as
  \texttt{jutari} (Julia) $\cdot$ \texttt{jaxtari} (JAX), gives four capabilities no neural
  network grants for free: exact forward, exact interventions $\mathrm{do}(u)$, full state
  observability, and content-path gradients. On it we build the oracle (centre): the primary
  intervention oracle of \Sectionname\ref{sec:oracle} (\Sectionname3.2), $\mathrm{do}(u{:=}v')\!\to\!\Delta y(u)$
  by bit-exact re-run, with a content-path gradient $\partial y/\partial u$ as a companion on
  the content path only. Any explanation $\hat{E}$ is scored by the correctness triad (right),
  $\text{right}=\mathrm{F}\wedge\mathrm{S}\wedge\mathrm{M}$, and reported on two axes,
  faithfulness against the oracle versus a plausibility proxy. Faithfulness aggregates
  correlation, sign and magnitude agreement, deletion/insertion area-under-curve, and
  precision at $k$ against the true causal top-$k$ (\Sectionname\ref{sec:triad}). A discrete index --- a
  sprite's position, set by strobe timing --- has a naive gradient that is zero under the hard
  forward map (Prop.~\ref{prop:zero}); for position, index, and event outputs the intervention
  oracle is the \emph{sole} ground truth and gradient methods are scored as methods under test.
  The schematic draws no experiment data. Data and oracle: see Supplement; oracle defined in
  \Sectionname\ref{sec:oracle} (\Sectionname3.2).}
  \label{fig:platform}
\end{figure}

\subsection{The differentiable, bit-exact VCS substrate}\label{sec:substrate}

The subject of this study is the Atari Video Computer System (VCS) itself, not a learned
agent. It is a complex real artifact never designed to be interpretable, yet whose every causal
dependency we can read off exactly. Paper~1 built it, a bit-exact, differentiable port of the
1977 VCS validated bit-for-bit against the \texttt{xitari} C++ reference on all 64 games of
the Paper-1 roster, in random-access memory and rendered screen, within the validated
conformance window~\cite{maier2025vcs,xitari}. Paper~1 is a technical companion that supplies
the mathematical proofs and experimental validation of the constructions we use here; no other
reference system like it exists anywhere, and no prior study has had the opportunity to
examine interpretability methods against this level of ground truth. We use the port in two
interchangeable, mutually cross-checked realizations, \texttt{jutari} (Julia) and
\texttt{jaxtari} (JAX/Python).

Paper~1 provides three execution modes. The \textbf{hard} mode is the bit-exact integer
machine, identical to the \texttt{xitari} C++ reference at every RAM byte and every screen
pixel. The \textbf{SOFT-STE} (straight-through estimator) mode keeps the forward pass
bit-identical to the hard mode at any finite temperature while routing a surrogate gradient
along the content path during the backward pass, a formulation whose forward equivalence
Paper~1 proves theoretically and validates experimentally across all 64 games. The
\textbf{soft} (fully relaxed) mode relaxes every discrete decision --- opcode dispatch via
a temperature-parameterised softmax, conditional branches via a sigmoid with steepness
$\alpha$, memory reads via a distance-weighted soft read --- into a fully differentiable
form. Paper~1 proves that the soft forward pass reproduces the hard machine bit-for-bit in
the limit $T \to 0^+$, $\alpha \to \infty$, and provides the temperature-limit error bound
$\|\Phi_{\mathrm{H}} - \Phi_{\mathrm{S}}\| \le (K-1)e^{-\Delta_{\min}/T} +
e^{-\alpha m_{\min}} + \rho e^{-1/T}$. At practical settings ($T \lesssim 0.14$, $\alpha
\gtrsim 6$), the soft mode remains forward-exact within machine precision for the validated
frame window; beyond those settings the forward pass diverges and the gradient quality
degrades.

The substrate grants four capabilities that no neural network offers for free. The first is
an exact forward map,
\begin{equation}
  y \;=\; \mathrm{VCS}\big(\mathrm{ROM},\, \mathrm{inputs}[0{:}t],\, s_0\big),
  \label{eq:vcs}
\end{equation}
that is, the output $y$ at step $t$, a pixel, a score, or a game event, is a deterministic
function of the cartridge bytes, the inputs up to $t$, and the initial state $s_0$, and
re-running reproduces $y$ to the bit. The second is exact interventions. We set any cause
$u$, a ROM byte, a RAM cell, a CPU/TIA/RIOT register, an opcode, or a joystick input, to any
value and re-run Eq.~\eqref{eq:vcs} deterministically. This realizes the causal operator
$\mathrm{do}(u)$ on a real machine rather than an assumed world model. The third is full
state observability, every bit of RAM and every register recorded along the trajectory, so
the ``activations'' we interpret are the program's own variables. The fourth is a gradient
$\partial y / \partial u$ that exists in all three modes. In SOFT-STE mode the gradient
routes through the content path via the straight-through estimator, while in soft mode the
gradient flows through every discrete decision, including the index boundary, at a cost of
numerical precision controlled by $T$ and $\alpha$. The inputs we
drive, the state we read and clamp, and the outputs we explain are the three faces a neural
interpretability study works with, so one substrate carries all three batteries.

We honor one boundary throughout. Paper-1 bit-exactness holds within a bounded frame window
on a fixed action stream, on the order of 30 frames of RAM and 60 frames of screen, and
every scored claim is taken inside that validated horizon. Long-horizon end-to-end gradients
are not needed for our single-step metrics, so we do not assume them.

\subsection{The ground-truth oracle}\label{sec:oracle}

The oracle is the heart of the method: for any output $y$ and any candidate cause $u$ it
returns the true effect of $u$ on $y$, against which every method is then scored. Its
primary construction is the exact intervention oracle: we intervene on $u$ by occlusion,
clamping, replacement, or resampling, re-run the machine deterministically, and record the
change in the output,
\begin{equation}
  \Delta y(u) \;=\; \mathrm{VCS}\big(\dots,\, \mathrm{do}(u{:=}v'),\, \dots\big) \;-\; y .
  \label{eq:doy}
\end{equation}
The true causal effect of a cause is thus the bit-exact difference its surgical edit makes
to the output, measured by actually running the altered machine, with no world model or
smoothness assumed. The collection $\{\Delta y(u)\}$ over all causes is the true causal map
of $y$, the reference for the faithfulness axis of every method.

The output $y$ comes in five regimes that fix what a method must predict and whether a
gradient is even meaningful. A \emph{content} output is a register, color, or graphics-bit
value that flows smoothly into the pixel, where the naive gradient is valid. A
\emph{position/index} output is a sprite, missile, or ball coordinate set by strobe timing,
where the naive gradient vanishes by Prop.~\ref{prop:zero} and only the bilinear-sampler
surrogate or the intervention oracle applies. An \emph{event} output is a discrete state
change, a collision latch or a life lost, also gradient-zero. A \emph{score} output is the
game's running counter, valid for the gradient only where it varies continuously. A
\emph{pixel} output is the rendered framebuffer value, content-valid where it tracks a smooth
cause and index-zero where it tracks a position. The Supplement lists, per game, which outputs
each method is scored on and the per-output gradient verdict.

The interventions that build the oracle are of distinct manifold kinds, and we state each
because the kind bounds what a score means. Occlusion and clamping are surgical but
\emph{off-manifold}: they can set a byte to a value the running program would never write.
Replacement from a donor trajectory is \emph{on-manifold}: the substituted value is one the
machine actually produced elsewhere. Resampling draws the cause from its own marginal over the
trajectory, and counterfactual edits are valid under the emulator by construction, since any
byte assignment yields a well-defined deterministic re-run. Off-manifold interventions are
acceptable for oracle scoring precisely because the oracle measures a \emph{causal} effect, not
a likelihood. The question is what the output would be if $u$ were $v'$. The bit-exact
machine answers it exactly for any $v'$, on- or off-manifold. Where a method's own claim is
about typical behavior rather than causal dependence we fall back to the on-manifold donor and
resampling variants and flag the choice in the Supplement.

The gradient $\partial y / \partial u$ through the differentiable substrate is a companion
construction. Its behaviour depends on the execution mode. In \textbf{SOFT-STE} mode, the
straight-through gradient routes to the content path only: register, colour, and graphics-bit
values that flow smoothly into the pixel agree with the intervention oracle, but a discrete
index --- a sprite's screen position set by the timing of a strobe write --- has a naive
gradient that vanishes under the hard forward map. In \textbf{soft} mode, the relaxation of
every discrete decision allows the gradient to flow through the index boundary as well, at a
cost controlled by the temperature $T$ and branch steepness $\alpha$. At the recommended
operating point ($T \approx 0.14$, $\alpha \gtrsim 6$), the soft gradient provides a
non-zero signal on position outputs while the forward pass remains bit-exact. As $T$ grows
or $\alpha$ shrinks, the forward pass diverges and the gradient becomes a looser surrogate.
We make the hard-map limit precise.

\begin{proposition}[Zero naive gradient for discrete-index outputs]\label{prop:zero}
Let an output $y$ depend on a cause $u$ only through a discrete index
$n(u)=\lfloor g(u)\rfloor\in\mathbb{Z}$, the integer count of colour clocks between a strobe
write and the position-counter reset that the VCS uses to place a sprite, missile, or ball,
with $g$ continuous, piecewise linear, and non-constant (strictly monotone on each linear
piece). Under the hard (bit-exact) forward map the rendered output is $y=h(n(u))$ for a fixed
render function $h$. Then $n(\cdot)$ is locally constant wherever $g(u)\notin\mathbb{Z}$, so
$\partial y/\partial u=0$ at every such point, and $y$ is discontinuous on the set
$\{u:g(u)\in\mathbb{Z}\}$, which is measure zero precisely because $g$ is strictly monotone on
each piece, and there the gradient is undefined. Hence the naive gradient is zero wherever the
derivative exists: an infinitesimal change in $u$ cannot move the sprite by a fractional
pixel. Any method that reduces to this naive gradient inherits the failure on position, index,
and event outputs.
\end{proposition}

The proposition states only the hard-map limit; it does not say that no useful gradient
exists. Paper~1~\cite{maier2025vcs} provides three constructions that escape it, each
corresponding to an execution mode. The \textbf{SOFT-STE} mode, through the straight-through
estimator, keeps the forward pass bit-exact at any finite temperature while routing a
non-zero gradient along the content path~\cite{bengio2013estimating,jang2017gumbel}; on
position outputs this gradient remains content-only and therefore zero. The \textbf{soft}
mode relaxes the discrete index boundary via a temperature-parameterised softmax, allowing
the gradient to flow through the index itself; the bilinear sampler is a special case of
this relaxation for spatial coordinates~\cite{jaderberg2015spatial}. The \textbf{bilinear
sampler} interpolates the index boundary to recover a usable position gradient without
relaxing the full machine~\cite{jaderberg2015spatial}. Paper~1 constructs and validates all
three, and proves the forward equivalence and temperature-limit bounds we assume.

Hence we have three companion gradient constructions, not one. The intervention oracle of
Eq.~\eqref{eq:doy} remains the primary ground truth across all modes. The SOFT-STE gradient
agrees with it on the content path but vanishes on position outputs; the soft gradient can
cover position outputs at the cost of temperature-dependent precision; and the bilinear
sampler provides a position gradient without relaxing the full machine. We score every
gradient-based method under all applicable modes and report the results separately. Where
the modes agree we report the common value; where they disagree the divergence itself is
informative about the method's reliance on gradient quality.

\subsection{Tiered ground truth and how we obtain it}\label{sec:tiers}

The ground truth comes in three tiers, and the tiering is not bookkeeping; it carries the
central argument of the paper. Tier~1 (T1) is the exact causal map of Eq.~\eqref{eq:doy},
derived internally from the substrate and available for all 64 games. Tier~2 (T2) is the
hardware roles and the read/write data-flow graph, also exact and internal across all 64
games. T1 and T2 are everything the machine \emph{is} and need no external label.

Tier~3 (T3) is different in kind, and that difference is the point. It attaches a
game-concept name to a cell or feature, ``the ball's $x$ position'' or ``the score,'' names
not computable from the substrate alone, because they are facts about what the program
\emph{means}, not what it \emph{does}. We import T3 from two public sources and make it
causal. AtariARI provides labels for 22 games, read from commented
disassemblies, with raw RAM coordinates not aligned to the rendered
position~\cite{anand2019unsupervised}; OCAtari covers more than 40 games and, in its RAM
mode, applies the render-time offsets that fix that misalignment~\cite{delfosse2023ocatari}.
We import both as candidates and verify each causally by setting the byte, re-rendering, and
confirming on the bit-exact framebuffer that the named object moves to the predicted place,
which upgrades a correlational label to a verified-causal one and auto-corrects the offsets;
further labels we discover by RAM-to-framebuffer correlation and intervention. The headline
core set is six games, Pong, Breakout, Space Invaders, Seaquest, Ms.\ Pac-Man, and Q*bert,
each carrying both sources, all in the Paper-1 bit-exact roster, together spanning
paddle-and-ball, shooter-with-shields, resource-meter, multi-agent-pursuit, and discrete-grid
logic. The remaining games form a breadth set scored on T1 and T2 only.

We are deliberate about what this buys us. All T3 labels are external and only verified,
never discovered from scratch, so a method that matches a supplied concept has aligned to a
given meaning, not recovered it. This is exactly why a probe can show that information is
\emph{present} without showing it is \emph{used}~\cite{hewitt2019designing}, and why our
intervention test, which asks whether the byte \emph{causes} the named object to move, is
strictly stronger than decoding it. Recovering the program's documentation and design from
the bare machine is a separate problem, the subject of the companion paper.

\subsection{The F\,$\wedge$\,S\,$\wedge$\,M correctness triad}\label{sec:triad}

A faithful attribution is not yet an understanding, and to make that a measurement we need an
operational criterion. We adopt one in the spirit of Marr's levels of
analysis~\cite{marr1982vision} and of Jonas and Kording's warning that finding structure is
not the same as explaining it~\cite{jonas2017could}: each axis is first a continuous score in
$[0,1]$, and the boolean ``right'' is a thresholded reading of those scores. We score an
explanation $\hat{E}$ of how an output $y$ arises by three numbers,
$\mathrm{F}_{\mathrm{score}}(\hat{E}),\,\mathrm{S}_{\mathrm{score}}(\hat{E}),\,
\mathrm{M}_{\mathrm{score}}(\hat{E})\in[0,1]$, and call it \emph{right} when all three clear
their thresholds,
\begin{equation}
  \mathrm{right}(\hat{E}) \;=\; \mathbf{1}\!\left[\,
  \mathrm{F}_{\mathrm{score}}\!\ge\!\tau_{\mathrm{F}} \,\wedge\,
  \mathrm{S}_{\mathrm{score}}\!\ge\!\tau_{\mathrm{S}} \,\wedge\,
  \mathrm{M}_{\mathrm{score}}\!\ge\!\tau_{\mathrm{M}}\,\right],
  \label{eq:triad}
\end{equation}
with $\tau_{\mathrm{F}}=\tau_{\mathrm{S}}=\tau_{\mathrm{M}}=0.5$ throughout. Results report the
continuous scores; the predicate is a derived summary, not the primary quantity. All three are
defined against the oracle, not against any method under test, which keeps the criterion
non-circular.

Faithfulness asks whether the claimed causes are the \emph{true} causes, scored against the
intervention oracle of Eq.~\eqref{eq:doy}. The per-record metric depends on the explanation's
output type, so $\mathrm{F}_{\mathrm{score}}$ is a single number assembled from three
type-matched primitives. The committed score is the bare effect-agreement primitive, not a
rescaled one: for a real-valued attribution map we report the Pearson (or Spearman, for ranked
outputs) correlation $\rho(\hat{E},\{\Delta y(u)\})$ between the attribution and the true causal
effects, scored \emph{raw}, so that a method whose attribution carries no causal signal lands at
$\rho=0$ rather than at the $0.5$ of an affine $\tfrac{1}{2}(1+\rho)$ remap. This is what makes
the position regime informative: under Prop.~\ref{prop:zero} the naive gradient is identically
zero there, its correlation with the oracle is $0$, and the reported faithfulness is therefore
$0.00$ and not a $0.5$ floor. The two committed estimators differ only in how the raw
correlation is summarized across the six core games. Phase~B reports the raw correlation
directly (e.g.\ occlusion at $\rho=0.482$), so a negative correlation is recorded as a negative
number; Phase~A's single-unit and tuning analyses, where per-game correlations can be negative,
report the per-game correlation clipped at $0$ and then averaged across games (e.g.\ the tuning
analysis A3, lowest in Phase~A at $0.116$, the clip-at-zero mean of six per-game correlations
that include negatives). For the same attribution map we also report, as quality measures,
deletion/insertion area-under-curve and precision at $k$ against the true causal top-$k$,
$\mathrm{P@}k=|\,\mathrm{top}_k(\hat{E})\cap\mathrm{top}_k(\{\Delta y(u)\})\,|/k$. For a
graph-valued explanation (Phase~A connectomics, Phase~C circuit discovery) we use the edge
$F_1$ against the true data-flow graph, the harmonic mean of precision and recall over recovered
edges, and for Phase~C patching and mediation the causal-effect-agreement correlation against
the true $\{\Delta y(u)\}$. The final $\mathrm{F}_{\mathrm{score}}$ is the per-record primary
metric averaged first within a game over its output records, then with equal weight across the
six core games, so no game-rich method dominates; confidence intervals are 95\% bootstrap over
games (\Sectionname\ref{sec:scoring}). The three Phases differ only in which primitive is primary: Phase~A
uses graph $F_1$ for the data-flow methods and the clip-at-zero mean correlation for the
single-unit and tuning analyses, Phase~B uses the raw effect-agreement correlation $\rho$ with
deletion/insertion AUC and $\mathrm{P@}k$ as the attribution map's quality measures, and Phase~C
uses the causal-effect-agreement correlation for patching and mediation and graph $F_1$ for
circuit discovery. The leaderboard score is therefore a single scalar per method, but it is
multi-metric in origin, and the per-method choice is recorded in the Supplement.

Sufficiency asks whether the explanation could regenerate the behavior under interventions it
has not seen. We treat $\mathrm{S}_{\mathrm{score}}$ as a held-out predictive estimator.
$\hat{E}$ is fit (or its causal sites read off) on a calibration set of interventions. It is
then asked to predict the true output $y'$ under a disjoint test set of held-out $\mathrm{do}(u)$
edits, with the prediction compared to the bit-exact re-run. The score is the fraction of
held-out interventions whose predicted $y'$ matches the true $y'$ within a fixed
output-type tolerance. The match is exact for discrete indices and events, and within a small
$\varepsilon$ band for continuous scores. Hence $\mathrm{S}_{\mathrm{score}}\in[0,1]$; a value
near zero means the explanation predicts no better than the unperturbed output. We deliberately
do not define it through interchange interventions or scrubbing. Those are themselves
mechanistic-interpretability methods under test in Phase~C, and using them would make the
criterion circular. Only Phase~B and Phase~C report $\mathrm{S}$. Its bootstrap intervals
follow \Sectionname\ref{sec:scoring}.

Minimality asks whether the explanation is parsimonious and lives at the algorithmic level of
the machine's own module decomposition. We define
$\mathrm{M}_{\mathrm{score}}=|U^\star|/|U_{\hat{E}}|\in(0,1]$, the ratio of the size of the
true minimal cause set $U^\star$, the smallest set of registers, RAM cells, opcodes, or
program variables whose intervention reproduces $y$, to the size $|U_{\hat{E}}|$ of the cause
set the explanation actually names. An explanation that names exactly the minimal set scores
$1$; one that records the whole machine state pays for every spurious cell. The target level
is the algorithmic unit, the register, RAM cell, opcode, module, or program variable, not the
pixel. This is why whole-state recording, which names all $128$ RAM cells where the true cause
set averages a handful, scores $\mathrm{M}_{\mathrm{score}}=0.07$ across the six core games:
its faithfulness is perfect because it withholds nothing, but it is maximally non-minimal.

Two reporting axes carry the headline comparison, faithfulness against the ground truth and a
plausibility proxy. The plausibility proxy is a documented, tradition-level value, not a
human-subjects measurement, and we report it as such. It is assigned per method tradition, not
per method instance, from the published evidence on how convincing each tradition's
explanations are taken to be, before any faithfulness score on our machine is computed, so the
proxy cannot be tuned to the result; the per-tradition assignment, its rubric, and its
sources are tabulated in the Supplement (Sec.~\ref{supp:plausibility}), where a sensitivity
analysis perturbs the proxy under additive, rank-preserving, rank-jittering, and
leave-one-tradition-out schemes and finds the faithfulness-versus-proxy rank correlation
stays negative throughout (from $-0.79$ to $-0.40$), so the headline contrast survives every
alternative proxy assignment we tried. Its job is to expose the danger zone of explanations with high
plausibility but low faithfulness, convincing yet wrong. Separating F, which the wiring
satisfies, from S and M, which the \emph{computation} must satisfy, is what lets us name,
later, the residual that faithfulness alone leaves behind.

\subsection{Methods under test, games, and aggregation}\label{sec:scoring}

We turn three traditions onto the machine and score each by the triad of
Eq.~\eqref{eq:triad}. The neuroscience battery (Phase~A) runs connectomics and data-flow
recovery, single-unit lesions, tuning curves, pairwise and global correlation structure,
pooled-activity spectra, Granger causality, dimensionality reduction, and whole-state
recording, scoring each against the known mechanism to quantify Kording's lesson at the
low-faithfulness baseline~\cite{jonas2017could}. The attribution and XAI
battery (Phase~B) runs gradient saliency and its variants, integrated gradients, occlusion,
perturbation masks, randomized masking, LIME, and Shapley-style
attribution~\cite{simonyan2014deep,sundararajan2017axiomatic,zeiler2014visualizing,fong2017interpretable,petsiuk2018rise,ribeiro2016why,lundberg2017unified};
Grad-CAM and its variants and attention
rollout~\cite{selvaraju2017gradcam,chattopadhyay2018gradcampp,abnar2020quantifying} require
neural layers or a learned policy and are recorded as not applicable. The
mechanistic-interpretability battery (Phase~C) runs activation patching and causal
mediation, interchange interventions and distributed alignment search, attribution and path
patching, automatic circuit discovery, sparse autoencoders, causal scrubbing, linear
probing with control tasks, and the logit and tuned
lens~\cite{vig2020investigating,meng2022locating,geiger2021causal,geiger2023finding,nanda2023attribution,wang2022interpretability,conmy2023towards,cunningham2023sparse,chan2022causal,alain2017understanding,hewitt2019designing,belrose2023eliciting}.

All faithfulness scoring runs on the six core games of Section~\ref{sec:tiers}, where the
oracle's re-runs are exact. Cell- and pixel-level scores, correlation, deletion and insertion
area-under-curve, and precision at $k$, are taken against T1 and need no T3. Only object- and
concept-level scores draw on the external T3 labels, and we flag every such score as
alignment-to-given, not discovery. We aggregate the committed per-game records into one row
per method on the shared faithfulness-versus-plausibility axes of the Results leaderboard.
The headline contrast is the position regime, where the naive gradient is zero under the hard
forward map by Prop.~\ref{prop:zero}.
There, causal and intervention methods average a faithfulness of $0.41$ against $0.07$ for
gradient and correlational methods, a gap of $0.34$. The per-method extreme places an
exact causal method, activation patching, at faithfulness $1.000$ and the field's default
saliency at $0.000$ on all six games. The popular method nonetheless carries the higher
plausibility proxy ($0.9$ versus $0.5$). Against these numbers the semantic gap is measured.

% =====================================================================
% BIBTEX-NEEDED: full entries for the keys cited above that are NOT yet
% in references.bib. The keys already present in references.bib that this
% section reuses are: xitari, jonas2017could, sundararajan2017axiomatic,
% selvaraju2017gradcam, chattopadhyay2018gradcampp, jaderberg2015spatial,
% jang2017gumbel, bengio2013estimating, delfosse2023ocatari.
% Please MERGE the following NEW entries (SM-only edit to references.bib):
% =====================================================================
%
% @article{maier2025vcs,
%   title   = {A bit-exact, differentiable port of the Atari 2600 Video Computer System},
%   author  = {Maier, Andreas and Bayer, Stefan and Krauss, Patrick},
%   journal = {(Paper 1; manuscript / preprint)},
%   year    = {2025},
%   note    = {jutari and jaxtari: 64/64 games bit-exact vs the xitari reference in RAM and screen}
% }
%
% @inproceedings{anand2019unsupervised,
%   title     = {Unsupervised State Representation Learning in Atari},
%   author    = {Anand, Ankesh and Racah, Evan and Ozair, Sherjil and Bengio, Yoshua and C{\^o}t{\'e}, Marc-Alexandre and Hjelm, R. Devon},
%   booktitle = {Advances in Neural Information Processing Systems (NeurIPS)},
%   year      = {2019}
% }
%
% @book{marr1982vision,
%   title     = {Vision: A Computational Investigation into the Human Representation and Processing of Visual Information},
%   author    = {Marr, David},
%   year      = {1982},
%   publisher = {W. H. Freeman and Company}
% }
%
% @inproceedings{simonyan2014deep,
%   title     = {Deep Inside Convolutional Networks: Visualising Image Classification Models and Saliency Maps},
%   author    = {Simonyan, Karen and Vedaldi, Andrea and Zisserman, Andrew},
%   booktitle = {International Conference on Learning Representations (ICLR) Workshop},
%   year      = {2014}
% }
%
% @inproceedings{zeiler2014visualizing,
%   title     = {Visualizing and Understanding Convolutional Networks},
%   author    = {Zeiler, Matthew D. and Fergus, Rob},
%   booktitle = {European Conference on Computer Vision (ECCV)},
%   year      = {2014},
%   pages     = {818--833}
% }
%
% @inproceedings{fong2017interpretable,
%   title     = {Interpretable Explanations of Black Boxes by Meaningful Perturbation},
%   author    = {Fong, Ruth C. and Vedaldi, Andrea},
%   booktitle = {IEEE International Conference on Computer Vision (ICCV)},
%   year      = {2017},
%   pages     = {3429--3437}
% }
%
% @inproceedings{petsiuk2018rise,
%   title     = {{RISE}: Randomized Input Sampling for Explanation of Black-box Models},
%   author    = {Petsiuk, Vitali and Das, Abir and Saenko, Kate},
%   booktitle = {British Machine Vision Conference (BMVC)},
%   year      = {2018}
% }
%
% @inproceedings{ribeiro2016why,
%   title     = {``Why Should I Trust You?'': Explaining the Predictions of Any Classifier},
%   author    = {Ribeiro, Marco Tulio and Singh, Sameer and Guestrin, Carlos},
%   booktitle = {Proceedings of the 22nd ACM SIGKDD International Conference on Knowledge Discovery and Data Mining},
%   year      = {2016},
%   pages     = {1135--1144}
% }
%
% @inproceedings{lundberg2017unified,
%   title     = {A Unified Approach to Interpreting Model Predictions},
%   author    = {Lundberg, Scott M. and Lee, Su-In},
%   booktitle = {Advances in Neural Information Processing Systems (NeurIPS)},
%   year      = {2017}
% }
%
% @inproceedings{abnar2020quantifying,
%   title     = {Quantifying Attention Flow in Transformers},
%   author    = {Abnar, Samira and Zuidema, Willem},
%   booktitle = {Proceedings of the 58th Annual Meeting of the Association for Computational Linguistics (ACL)},
%   year      = {2020},
%   pages     = {4190--4197}
% }
%
% @inproceedings{vig2020investigating,
%   title     = {Investigating Gender Bias in Language Models Using Causal Mediation Analysis},
%   author    = {Vig, Jesse and Gehrmann, Sebastian and Belinkov, Yonatan and Qian, Sharon and Nevo, Daniel and Singer, Yaron and Shieber, Stuart},
%   booktitle = {Advances in Neural Information Processing Systems (NeurIPS)},
%   year      = {2020}
% }
%
% @inproceedings{meng2022locating,
%   title     = {Locating and Editing Factual Associations in {GPT}},
%   author    = {Meng, Kevin and Bau, David and Andonian, Alex and Belinkov, Yonatan},
%   booktitle = {Advances in Neural Information Processing Systems (NeurIPS)},
%   year      = {2022}
% }
%
% @inproceedings{geiger2021causal,
%   title     = {Causal Abstractions of Neural Networks},
%   author    = {Geiger, Atticus and Lu, Hanson and Icard, Thomas and Potts, Christopher},
%   booktitle = {Advances in Neural Information Processing Systems (NeurIPS)},
%   year      = {2021}
% }
%
% @inproceedings{geiger2023finding,
%   title     = {Finding Alignments Between Interpretable Causal Variables and Distributed Neural Representations},
%   author    = {Geiger, Atticus and Wu, Zhengxuan and Potts, Christopher and Icard, Thomas and Goodman, Noah D.},
%   booktitle = {Proceedings of the Third Conference on Causal Learning and Reasoning (CLeaR)},
%   year      = {2024}
% }
%
% @misc{nanda2023attribution,
%   title        = {Attribution Patching: Activation Patching at Industrial Scale},
%   author       = {Nanda, Neel},
%   year         = {2023},
%   howpublished = {\url{https://www.neelnanda.io/mechanistic-interpretability/attribution-patching}}
% }
%
% @inproceedings{wang2022interpretability,
%   title     = {Interpretability in the Wild: A Circuit for Indirect Object Identification in {GPT-2} Small},
%   author    = {Wang, Kevin and Variengien, Alexandre and Conmy, Arthur and Shlegeris, Buck and Steinhardt, Jacob},
%   booktitle = {International Conference on Learning Representations (ICLR)},
%   year      = {2023}
% }
%
% @inproceedings{conmy2023towards,
%   title     = {Towards Automated Circuit Discovery for Mechanistic Interpretability},
%   author    = {Conmy, Arthur and Mavor-Parker, Augustine N. and Lynch, Aengus and Heimersheim, Stefan and Garriga-Alonso, Adri{\`a}},
%   booktitle = {Advances in Neural Information Processing Systems (NeurIPS)},
%   year      = {2023}
% }
%
% @misc{cunningham2023sparse,
%   title         = {Sparse Autoencoders Find Highly Interpretable Features in Language Models},
%   author        = {Cunningham, Hoagy and Ewart, Aidan and Riggs, Logan and Huben, Robert and Sharkey, Lee},
%   year          = {2023},
%   eprint        = {2309.08600},
%   archivePrefix = {arXiv},
%   primaryClass  = {cs.LG}
% }
%
% @misc{chan2022causal,
%   title        = {Causal Scrubbing: A Method for Rigorously Testing Interpretability Hypotheses},
%   author       = {Chan, Lawrence and Garriga-Alonso, Adri{\`a} and Goldowsky-Dill, Nicholas and Greenblatt, Ryan and Nitishinskaya, Jenny and Radhakrishnan, Ansh and Shlegeris, Buck and Thomas, Nate},
%   year         = {2022},
%   howpublished = {AI Alignment Forum}
% }
%
% @inproceedings{alain2017understanding,
%   title     = {Understanding Intermediate Layers Using Linear Classifier Probes},
%   author    = {Alain, Guillaume and Bengio, Yoshua},
%   booktitle = {International Conference on Learning Representations (ICLR) Workshop},
%   year      = {2017}
% }
%
% @inproceedings{hewitt2019designing,
%   title     = {Designing and Interpreting Probes with Control Tasks},
%   author    = {Hewitt, John and Liang, Percy},
%   booktitle = {Proceedings of the 2019 Conference on Empirical Methods in Natural Language Processing (EMNLP)},
%   year      = {2019},
%   pages     = {2733--2743}
% }
%
% @misc{belrose2023eliciting,
%   title         = {Eliciting Latent Predictions from Transformers with the Tuned Lens},
%   author        = {Belrose, Nora and Furman, Zach and Smith, Logan and Halawi, Danny and Ostrovsky, Igor and McKinney, Lev and Biderman, Stella and Steinhardt, Jacob},
%   year          = {2023},
%   eprint        = {2303.08112},
%   archivePrefix = {arXiv},
%   primaryClass  = {cs.LG}
% }
